\section{\HMTLang{Réciprocité différentielle et restitution des canaux constitutifs}{Réciprocité différentielle et récupération des canaux constitutifs}}
\label{delta22:iii:restitucion}
\HMTLang{Les canaux du vide, la fonctionnelle de Barbero--Immirzi et le registre dodécaphasique appartiennent à une même chaîne de constructions. Le résultat suivant réunit leurs inverses : la vitesse normalisée conserve la composante commune de la réponse, tandis que la lecture orientée de Barbero détermine sa répartition entre feuillets. La composition retrouve la paire angulaire et, sur l’image engendrée avec sa section dimensionnelle, le coefficient d’action. En conservant en outre les coordonnées régionales, elle atteint la lecture périodique de \(K\).}{Les canaux du vide, la fonctionnelle de Barbero--Immirzi et le registre dodécaphasique appartiennent à une même chaîne de constructions. Le résultat suivant réunit leurs inverses : la vitesse normalisée conserve la composante commune de la réponse, tandis que la lecture orientée de Barbero détermine sa distribution entre feuillets. La composition retrouve la paire angulaire et, sur l’image engendrée avec sa section dimensionnelle, le coefficient d’action. Conserver les coordonnées régionales donne également la lecture périodique de \(K\).}

\HMTLang{La construction précédente d’APP--TRIT--TPK et de son état enrichi demeure l’antécédent. Les paramètres employés ci-après sont les représentations angulaires déjà obtenues ; les preuves analytiques décrivent leur réalisation et les opérations de récupération. Les projecteurs étiquetés \(P_\pm\), \(\mathcal R=P_+-P_-\), l’orientation et le relèvement angulaire restent explicites.}{La construction précédente d’APP--TRIT--TPK et de son état enrichi demeure l’antécédent. Les paramètres utilisés ci-après sont les sorties angulaires déjà obtenues ; les preuves analytiques décrivent leur réalisation et leurs opérations de récupération. Les projecteurs étiquetés \(P_\pm\), \(\mathcal R=P_+-P_-\), l’orientation et le relèvement angulaire restent explicites.}

\subsection{\HMTLang{Potentiel spectral et réciprocité constitutive}{Potentiel spectral et réciprocité constitutive}}
\HMTLang{Sur le cône \(x>|y|\), soient}{Sur le cône \(x>|y|\), soient}
\[
 T=e^{-xI-y\mathcal R},\quad s_\pm=e^{-30(x\pm y)},\quad
 f(s)=\frac{1-s^3}{1-s^4},\quad
 V=(I-T^{90})(I-T^{120})^{-1}=r_+P_++r_-P_-,
\]
\[
 r_\pm=f(s_\pm),\quad
 \widehat c=(r_+r_-)^{-1},\quad \widehat Z=r_-/r_+,
 \quad c_\ell=\log\widehat c,\quad z_\ell=\log\widehat Z.
\]
\HMTLang{L’inversibilité de \(I-T^{120}\) découle de \(0<s_\pm<1\). Le quotient de polynômes prolongé est \(f(s)=(1+s+s^2)/(1+s+s^2+s^3)\) ; ses extrémités sont \(1\) et \(3/4\).}{L’inversibilité de \(I-T^{120}\) découle de \(0<s_\pm<1\). Le quotient de polynômes prolongé est \(f(s)=(1+s+s^2)/(1+s+s^2+s^3)\) ; ses valeurs aux extrémités sont \(1\) et \(3/4\).}
\begin{theorem}[\HMTLang{Potentiel de la réponse normalisée}{Potentiel de la réponse normalisée}]
\HMTLang{La série}{La série}
\begin{equation}
 \mathscr F(x,y)=\sum_{\sigma=\pm1}\sum_{n\ge1}
 \left[\frac{e^{-120n(x+\sigma y)}}{120n^2}
       -\frac{e^{-90n(x+\sigma y)}}{90n^2}\right]
 \label{delta22:iii:potencial}
\end{equation}
\HMTLang{définit un potentiel analytique réel strictement concave, avec}{définit un potentiel analytique réel strictement concave, avec}
\[
 \nabla\mathscr F=(c_\ell,z_\ell),\qquad
 \partial_y\log\widehat c=\partial_x\log\widehat Z.
\]
\HMTLang{La normalisation est \(\mathscr F\to0\) lorsque \(x-|y|\to\infty\).}{Sa normalisation est \(\mathscr F\to0\) lorsque \(x-|y|\to\infty\).}
\end{theorem}
\begin{proof}
\HMTLang{Sur chaque compact du cône, il existe \(\epsilon>0\) avec \(x\pm y\ge\epsilon\). Chaque dérivée d’ordre fixé des termes est dominée par une puissance de \(n\) multipliée par \(e^{-90n\epsilon}\) ; la somme converge uniformément. La dérivation terme à terme et \(\sum_{n\ge1}z^n/n=-\log(1-z)\) donnent les deux composantes indiquées. Pour la concavité, soient}{Sur chaque partie compacte du cône, il existe \(\epsilon>0\) avec \(x\pm y\ge\epsilon\). Toute dérivée d’ordre fixé d’un terme est bornée par une puissance de \(n\) multipliée par \(e^{-90n\epsilon}\) ; la somme converge uniformément. La dérivation terme à terme et \(\sum_{n\ge1}z^n/n=-\log(1-z)\) donnent les composantes indiquées. Pour la concavité, posons}
\[
 g(w)=\log f(e^{-w}),\quad u=30(x+y),\quad v=30(x-y),
 \quad a=g'(u),\quad b=g'(v).
\]
\[
 g'(w)=\frac3{e^{3w}-1}-\frac4{e^{4w}-1}>0,
 \qquad
 D(c_\ell,z_\ell)=-30
 \begin{pmatrix}a+b&a-b\\a-b&a+b\end{pmatrix}.
\]
\HMTLang{L’inégalité s’obtient parce que \(t/(e^{tw}-1)\) décroît strictement en \(t>0\) : sa dérivée a pour numérateur \(e^{tw}(1-tw)-1<0\). Les valeurs propres de la matrice sont \(-60a\) et \(-60b\), toutes deux négatives, et son déterminant est \(3600ab>0\). La symétrie prouve la réciprocité différentielle. La borne exponentielle donne la limite de normalisation.}{L’inégalité découle du fait que \(t/(e^{tw}-1)\) est strictement décroissante pour \(t>0\) : sa dérivée a pour numérateur \(e^{tw}(1-tw)-1<0\). Les valeurs propres de la matrice sont \(-60a\) et \(-60b\), toutes deux négatives, et son déterminant est \(3600ab>0\). La symétrie prouve la réciprocité différentielle. La borne exponentielle donne la limite de normalisation.}
\end{proof}
\HMTLang{Le même moment spectral \(\mathcal D_2(z)=\sum_{n\ge1}z^n/n^2\) permet d’écrire \eqref{delta22:iii:potencial} au moyen de ses évaluations réelles contractantes. Le caractère orienté déjà construit satisfait \(\mathcal C_2(s)=\operatorname{Im}\mathcal D_2(is)\), \(\mathcal C_2(1)=G_{\rm Cat}\). La convergence absolue atteint \(s=1\) ; pour les séries dérivées qui cessent d’y converger, on conserve la limite radiale d’Abel. Les évaluations réelles et orientées partagent l’indice de profondeur et la pondération, en conservant leurs arguments respectifs. Sous l’amplification \(T\otimes I_{9^m}\), la trace divisée par \(9^m\) conserve ces lectures, car chaque valeur propre se répète exactement \(9^m\) fois. Cette identité correspond à cette amplification déclarée.}{Le même moment spectral \(\mathcal D_2(z)=\sum_{n\ge1}z^n/n^2\) exprime \eqref{delta22:iii:potencial} au moyen d’évaluations réelles contractantes. Le caractère orienté déjà construit satisfait \(\mathcal C_2(s)=\operatorname{Im}\mathcal D_2(is)\), \(\mathcal C_2(1)=G_{\rm Cat}\). La convergence absolue s’étend jusqu’à \(s=1\) ; pour les séries dérivées qui cessent d’y converger, on conserve la limite radiale d’Abel. Les évaluations réelles et orientées partagent l’indice de profondeur et la pondération tout en conservant leurs arguments respectifs. Sous l’amplification par \(T\otimes I_{9^m}\), la trace divisée par \(9^m\) préserve ces lectures, puisque chaque valeur propre est répétée exactement \(9^m\) fois. Cette identité concerne l’amplification spécifiée.}

\begin{proposition}[\HMTLang{Sensibilité orientée et convexité}{Sensibilité orientée et convexité}]
\HMTLang{Pour \(x>0\) fixé, \(y\mapsto c_\ell(x,y)\) est paire et strictement convexe, avec un minimum unique en zéro.}{Pour \(x>0\) fixé, \(y\mapsto c_\ell(x,y)\) est paire et strictement convexe, avec son minimum unique en zéro.}
\end{proposition}
\begin{proof}
\[
 g''(w)=-\frac9{4\sinh^2(3w/2)}+\frac{16}{4\sinh^2(2w)}<0.
\]
\HMTLang{En effet, \(t/\sinh(tw/2)\) décroît : sa dérivée a le signe de \(\sinh v-v\cosh v<0\). Comme \(c_\ell=-g(u)-g(v)\), on obtient \(\partial_y^2c_\ell=-900[g''(u)+g''(v)]>0\). La symétrie sous \(y\mapsto-y\) fixe le minimum. Le développement local est}{En effet, \(t/\sinh(tw/2)\) décroît : sa dérivée a le signe de \(\sinh v-v\cosh v<0\). Puisque \(c_\ell=-g(u)-g(v)\), on obtient \(\partial_y^2c_\ell=-900[g''(u)+g''(v)]>0\). La symétrie sous \(y\mapsto-y\) fixe le minimum. Le développement local est}
\[
 c_\ell=-2g(30x)-900g''(30x)y^2+O(y^4),\qquad
 z_\ell=-60g'(30x)y+O(y^3).
\]
\end{proof}
\HMTLang{La vitesse retient la séparation des canaux au second ordre ; l’impédance conserve leur orientation au premier ordre. Évaluer le canal moyen avant le lecteur élimine précisément ces contributions.}{La vitesse retient la séparation des canaux au second ordre ; l’impédance retient leur orientation au premier ordre. Évaluer le canal moyen avant d’appliquer le lecteur élimine précisément ces contributions.}

\subsection{\HMTLang{Inversion globale au moyen de la vitesse et de Barbero--Immirzi}{Inversion globale au moyen de la vitesse et de Barbero--Immirzi}}
\HMTLang{L’inverse \(\psi=f^{-1}:(3/4,1)\to(0,1)\) est défini par l’unique racine de}{L’inverse \(\psi=f^{-1}:(3/4,1)\to(0,1)\) est défini par l’unique racine de}
\[
 rs^3+(r-1)(1+s+s^2)=0.
\]
\HMTLang{L’unicité découle de}{L’unicité découle de}
\[
 f'(s)=-\frac{s^2(s^2+2s+3)}{(1+s+s^2+s^3)^2}<0.
\]
\HMTLang{Nous conservons la fonctionnelle de retour déjà dérivée de l’incidence hexade--octade :}{Conservons la fonctionnelle de retour déjà dérivée de l’incidence hexade--octade :}
\[
 \mathcal H(s)=\frac{12s^3}{1-s^9}-\frac{s^4}{1-s^{12}}
              -\frac{(\log s)^2}{20\pi},\qquad F_{\rm B}=\mathcal H\circ\psi.
\]
\begin{theorem}[\HMTLang{Restitution constitutive orientée}{Récupération constitutive orientée}]
\HMTLang{L’application}{L’application}
\[
 (r_+,r_-)\longmapsto
 \left(\widehat c=\frac1{r_+r_-},\quad
 \gamma=F_{\rm B}(r_-)-F_{\rm B}(r_+)\right)
\]
\HMTLang{est un difféomorphisme analytique réel de \((3/4,1)^2\) sur \((1,16/9)\times\mathbb R\). Cette affirmation décrit le domaine analytique du lecteur ; l’image effectivement engendrée conserve, en outre, les compatibilités du TPK.}{est un difféomorphisme analytique réel de \((3/4,1)^2\) sur \((1,16/9)\times\mathbb R\). Cette affirmation décrit le domaine analytique du lecteur ; l’image effectivement engendrée conserve en outre les conditions de compatibilité du TPK.}
\end{theorem}
\begin{proof}
\[
 \mathcal H'(s)=\frac{36s^2(1+2s^9)}{(1-s^9)^2}
 -\frac{4s^3(1+2s^{12})}{(1-s^{12})^2}
 -\frac{\log s}{10\pi s}>0.
\]
\HMTLang{Le quotient du premier terme positif par la grandeur du second est}{Le rapport du premier terme positif à la grandeur du second est}
\[
 \frac9s\frac{1+2s^9}{1+2s^{12}}
 \left(\frac{1-s^{12}}{1-s^9}\right)^2>9,
\]
\HMTLang{et le dernier terme de la somme est également positif. Les extrémités sont \(\mathcal H(s)\sim-(\log s)^2/(20\pi)\) en zéro et \(\mathcal H(s)\sim5/[4(1-s)]\) en un. Par conséquent, \(F_{\rm B}'<0\), avec les limites \(+\infty\) en \(3/4\) et \(-\infty\) en \(1\).}{et le dernier terme de la somme est également positif. Les comportements asymptotiques aux extrémités sont \(\mathcal H(s)\sim-(\log s)^2/(20\pi)\) en zéro et \(\mathcal H(s)\sim5/[4(1-s)]\) en un. Donc \(F_{\rm B}'<0\), avec les limites \(+\infty\) en \(3/4\) et \(-\infty\) en \(1\).}

\HMTLang{Pour \(c=\widehat c\) fixé, soit \(z=\log\widehat Z\). Alors}{Pour \(c=\widehat c\) fixé, posons \(z=\log\widehat Z\). Alors}
\[
 r_+=c^{-1/2}e^{-z/2},\quad r_-=c^{-1/2}e^{z/2},\quad
 |z|<a(c)=\min\{\log c,\log(16/(9c))\}.
\]
\HMTLang{La fonction \(\mathcal G_c(z)=F_{\rm B}(r_-)-F_{\rm B}(r_+)\) satisfait}{La fonction \(\mathcal G_c(z)=F_{\rm B}(r_-)-F_{\rm B}(r_+)\) satisfait}
\[
 \mathcal G_c'(z)=\tfrac12
 [r_-F_{\rm B}'(r_-)+r_+F_{\rm B}'(r_+)]<0.
\]
\HMTLang{À l’extrémité supérieure, \(r_-\to1\), ou \(r_+\to3/4\), ou les deux, et \(\mathcal G_c\to-\infty\). À l’extrémité inférieure, elle tend vers \(+\infty\). Le théorème des valeurs intermédiaires et la monotonie donnent une unique solution pour chaque \(\gamma\). La dérivée non nulle fournit l’inverse analytique par le théorème des fonctions implicites. La description par le minimum de \(a(c)\) délimite seulement la frontière ; l’argument s’effectue sur son domaine ouvert.}{À l’extrémité supérieure, \(r_-\to1\), ou \(r_+\to3/4\), ou les deux, et \(\mathcal G_c\to-\infty\). À l’extrémité inférieure, elle tend vers \(+\infty\). Le théorème des valeurs intermédiaires et la monotonie donnent une unique solution pour chaque \(\gamma\). La dérivée non nulle donne l’inverse analytique par le théorème des fonctions implicites. Le minimum dans \(a(c)\) décrit seulement la frontière ; l’argument se déroule sur son domaine ouvert.}
\end{proof}

\subsection{\HMTLang{Récupération angulaire, action et registre périodique}{Récupération angulaire, action et registre périodique}}
\HMTLang{Une fois \(r_\pm\) obtenus, l’inversion cubique donne \(s_\pm\). Avec le relèvement et l’orientation conservés,}{Une fois \(r_\pm\) obtenus, l’inversion cubique donne \(s_\pm\). Avec le relèvement et l’orientation conservés,}
\[
 x=-\frac1{60}\log(s_+s_-),\quad y=\frac1{60}\log\frac{s_-}{s_+},
 \quad A_{\deg}=-\frac3\pi\log(s_+s_-),\quad
 C^*_{\deg}=\frac3\pi\log\frac{s_-}{s_+}.
\]
\HMTLang{Sur la section positive engendrée \(A_{\deg}=1000\alpha\), \(C^*_{\deg}=2(\eta_{\rm ret}+\alpha)\), on retrouve}{Sur la section positive engendrée \(A_{\deg}=1000\alpha\), \(C^*_{\deg}=2(\eta_{\rm ret}+\alpha)\), on retrouve}
\[
 \alpha=A_{\deg}/1000,\qquad
 \eta_{\rm ret}=C^*_{\deg}/2-\alpha,\qquad
 \hbar_{\rm ret}=s_0\eta_{\rm ret},\qquad s_0:=10^{-34}\mathcal U_S.
\]
\HMTLang{L’échelle \(s_0\) et les projecteurs sont des antécédents conservés. Dans l’orientation conjuguée, l’étiquette d’action est également transportée. La récupération est un inverse sur l’image correspondante.}{L’échelle \(s_0\) et les projecteurs sont des antécédents conservés. Sous l’orientation conjuguée, l’étiquette d’action est également transportée. La récupération est un inverse sur l’image correspondante.}

\begin{proposition}[\HMTLang{Composition avec la lecture périodique de douze blocs}{Composition avec la lecture périodique de douze blocs}]
\HMTLang{La lecture périodique prolonge les douze blocs de \(K\), tandis que les coordonnées régionales conservent leur génération coinductive complète. Avec \(\alpha_A\) comme coordonnée de la clôture arithmétique précédemment définie,}{La lecture périodique prolonge les douze blocs de \(K\), tandis que les coordonnées régionales conservent leur génération coinductive complète. Avec \(\alpha_A\) désignant la coordonnée de clôture arithmétique précédemment définie,}
\[
 \kappa_{\rm per}=\pi_{\rm HMT}+e_{\rm HMT}-\varphi_{\rm HMT}-4-\alpha_A,
 \quad B=1000,\quad N_K=(B^{12}-1)\kappa_{\rm per}.
\]
\HMTLang{Sur l’image du registre engendré, \(N_K\) est entier et ses douze blocs, avec origine fixe et zéros initiaux conservés, se retrouvent exactement par}{Sur l’image du registre engendré, \(N_K\) est un entier et ses douze blocs, avec origine fixe et zéros initiaux conservés, se retrouvent exactement par}
\[
 K_j=\left\lfloor\frac{N_K}{B^{12-j}}\right\rfloor
 -B\left\lfloor\frac{N_K}{B^{13-j}}\right\rfloor,
 \qquad 1\le j\le12.
\]
\end{proposition}
\begin{proof}
\HMTLang{La division euclidienne successive de \(N_K\) en base \(B\) détermine de manière unique \(K_j\in\{0,\ldots,B-1\}\) et produit \(N_K=\sum_{j=1}^{12}K_jB^{12-j}\). La répétition périodique de ces blocs somme la série géométrique \(N_K\sum_{m\ge1}B^{-12m}=N_K/(B^{12}-1)\). Les deux compositions sont inverses. La lecture finie \(N_K/B^{12}\) conserve une autre normalisation, distincte de la périodique. Avec des données approchées, la récupération exige un intervalle certifié pour \(N_K\) contenant un unique entier.}{La division euclidienne successive de \(N_K\) en base \(B\) détermine de manière unique \(K_j\in\{0,\ldots,B-1\}\) et donne \(N_K=\sum_{j=1}^{12}K_jB^{12-j}\). La répétition périodique de ces blocs somme la série géométrique \(N_K\sum_{m\ge1}B^{-12m}=N_K/(B^{12}-1)\). Les deux compositions sont inverses. La lecture finie \(N_K/B^{12}\) conserve une normalisation différente de celle de la lecture périodique. Avec des données approchées, la récupération nécessite un intervalle certifié pour \(N_K\) contenant un unique entier.}
\end{proof}
\HMTLang{Cette composition réunit la restitution constitutive et la récupération du registre déjà construit. La sélection exceptionnelle et temporelle détermine \(K\) au préalable ; l’inverse explique quelle information ses lecteurs permettent de restituer. La portée est le registre ordonné de douze coordonnées avec les antécédents indiqués.}{Cette composition réunit la restitution constitutive et la récupération du registre déjà construit. La sélection exceptionnelle et temporelle détermine \(K\) au préalable ; l’inverse explique quelle information ses lecteurs peuvent restituer. Sa portée est le registre ordonné de douze coordonnées avec les antécédents indiqués.}

\subsection{\HMTLang{Bornes de stabilité dans le domaine contractant}{Bornes de stabilité dans le domaine contractant}}
\HMTLang{Pour \(w>0\), écrivons \(a(w)=g'(w)\). Avec \(s=e^{-w}\),}{Pour \(w>0\), écrivons \(a(w)=g'(w)\). Avec \(s=e^{-w}\),}
\[
 a(w)=\frac{s^3(s^2+2s+3)}{(1+s+s^2)(1+s+s^2+s^3)},
 \qquad \tfrac12e^{-3w}<a(w)<3e^{-3w}.
\]
\HMTLang{Les bornes s’obtiennent en comparant les polynômes positifs du quotient ; \(g''<0\) démontre sa décroissance. Sur un domaine convexe où \(0<m_0\le30(x\pm y)\le M_0<\infty\), l’application \(\mathscr C=(c_\ell,z_\ell)\) satisfait, pour la norme euclidienne,}{Les bornes découlent de la comparaison des polynômes positifs du quotient ; \(g''<0\) prouve sa décroissance. Sur un domaine convexe où \(0<m_0\le30(x\pm y)\le M_0<\infty\), l’application \(\mathscr C=(c_\ell,z_\ell)\) satisfait, pour la norme euclidienne,}
\[
 60a(M_0)\|p-q\|\le\|\mathscr C(p)-\mathscr C(q)\|
 \le60a(m_0)\|p-q\|.
\]
\HMTLang{Pour la borne supérieure, on intègre la jacobienne sur le segment. Pour la borne inférieure, on applique le produit scalaire avec \(p-q\) à la jacobienne définie négative, puis Cauchy--Schwarz. La norme induite de la différentielle inverse est \(1/[60\min\{a(u),a(v)\}]\). En contraction profonde, elle peut croître exponentiellement, même lorsque le nombre de condition relatif reste proche de un. L’unicité exacte et la robustesse face à l’erreur sont donc spécifiées par des critères différents et compatibles.}{Pour la borne supérieure, on intègre la jacobienne le long du segment. Pour la borne inférieure, on prend le produit scalaire avec \(p-q\), on utilise le caractère défini négatif de la jacobienne, puis on applique Cauchy--Schwarz. La norme induite de la différentielle inverse est \(1/[60\min\{a(u),a(v)\}]\). En contraction profonde, elle peut croître exponentiellement même lorsque le nombre de condition relatif reste proche de un. L’unicité exacte et la robustesse face à l’erreur sont ainsi spécifiées par des critères distincts et compatibles.}
